\documentclass[11pt]{article}

\usepackage[margin=1in]{geometry}
\usepackage[T1]{fontenc}
\usepackage{lmodern}
\usepackage{microtype}
\usepackage{enumitem}
\usepackage[hidelinks]{hyperref}
\usepackage{xurl}
\usepackage{parskip}

\title{\textbf{20 Questions with LLMs}\\[0.5em]
\large CSE-519 Project Proposal}
\author{}
\date{}

\begin{document}

\maketitle

\section{Abstract (Introduction/General Plan)}

The game of 20 questions involves a subject (person) having an object,
profession or concept in their mind while players (guessers) question the
subject to narrow down and guess the object, having the limit of 20 questions
to do so. Our project aims to evaluate the performance of Large Language Models
(LLMs) for both the subject and the player.

Throughout the years, LLMs have gotten better at emulating human conversations
and at playing games at levels comparable (or even exceeding) professions. We
aim to see how well they have really gotten and which LLMs in particular
perform best at such tasks. Not only will we investigate their performance at
the game, but we will also judge the models on how human-like they sound when
posing questions. We will also look into reasons why certain LLMs are better
than others, investigating how they were trained and for which original purpose
they were built.

\section{Background (Previous Work Done in this Domain)}

According to C. R. S. Kumar in his paper \textit{Large Language Models (LLM)
for Automating 20 Questions Game}, the game used to be accomplished by ``using
decision trees, heuristic algorithms, or rule-based systems that map out
potential objects and the corresponding questions that might lead to their
identification. However, these methods often struggle with the complexity and
variability inherent in human language, especially when dealing with ambiguous
or less common objects.'' He goes on to utilize 1 LLM as a guesser, and another
LLM as the answerer. This paper then uses the data gathered from a Kaggle
competition, along with other sources.

In 2024, Kaggle hosted a competition to see who could create the best LLM for
the 20 Questions game. Each team had the ability to submit bots who would play
games with other bots, and the winner of the competition utilized a dual LLM
approach to handle different types of questions, with a Meta Llama 3 8B
Instruct as the general purpose LLM, and a DeepSeek Math model to handle
complex mathematical oriented questions about the keyword. The competition
featured a 2 vs 2 cooperative game, with each team consisting of a questioner
and answerer. Our work will be simplifying this to a one versus one
conventional implementation.

While previous work has demonstrated that LLMs can participate in 20 Questions,
there is still an opportunity to investigate how different language models
compare when assigned different roles.

An important consideration is that the questioner and answerer require
different capabilities. The questioner must demonstrate strategic reasoning,
effective question selection, and the ability to interpret previous responses.
The answerer must demonstrate factual knowledge, consistency, and the ability
to interpret questions accurately. Because of this, a model that performs well
as a questioner may not perform well as an answerer. Since our project requires
each model to perform both roles, the models chosen need to be well balanced
for both of these, and the winning model may not be the best guesser, or the
best answerer.

\section{Approach (Plan in Detail)}

Our approach involves finding LLMs and researching what they were originally
created for. Playing 20 questions directly with each LLM (as a
player/questioner) and evaluating performance based on how many questions it
takes for them to guess our object, and if they guessed correctly or not.

We will develop a scoring algorithm to score each LLM on its performance.

We will also investigate which topics each LLM performs best on.

We would like to evaluate the humanness of the questions the models come up
with for which we will also have to develop a metric. This might be subjective,
however, we will try our best to find characteristics that seem more bot-like,
rather than natural.

\textbf{Pairwise comparison:} We will put pairs of LLMs head to head (even with
themselves) and see how well they do against each other and what concepts they
come up with. Are the concepts repetitive?, which LLM does best when it comes
to coming up with concepts more randomly?

For the questions asked, we will evaluate and score them based on how good the
question was at narrowing the search down.

For each pair we will initially have a hypothesis as to how it will perform
based on how it was trained. Then we will see if our hypothesis holds true, and
what could explain their actual behaviour.

We will test multiple models and run brief experiments involving runs with a
small set of topics, then we will identify which ones work best and use them
for further analysis.

\section{Models and Data}

For many of the models we looked at, we were searching for four main things:

\begin{description}[leftmargin=2.2cm, style=nextline]
    \item[Knowledge:] How good it is with facts and scientific know-how.
    \item[Humanness:] How naturally it poses questions, as opposed to sounding
    like a bot.
    \item[Size:] We want to investigate how well small models perform compared
    to larger models. Does Occom's Razor hold? Do simpler/lightweight models do
    the job just as well?
    \item[Logic and Reasoning:] How well it does for logic and deduction, this
    would be key when it comes to narrowing down the questions to reach the
    answer faster.
\end{description}

\subsection*{Candidate Models}

\textbf{DeepSeek R1 0528:} DeepSeek is a powerful model when it comes to strong
reasoning and inference tasks. It performs well for several benchmarks
compared to models. Most notably the GPQA Diamond benchmark, which means it is
good with facts on major sciences (Physics, Biology and Chemistry). We believe
that this might give it an edge when it comes to 20 Questions, specially for
games that involve scientific facts to deduce the answer.

\textbf{Kimi K2:} Although originally meant for coding tasks, KIMI K2 is good
at conversations and step by step logical reasoning. We thought it would be
interesting to see how it's thinking style does for 20 questions, even though
it is op

\textbf{Xwinlm:} Xwin-LM is based on Llama 2 and is optimized for emulating
human-like conversations. I performed well on several benchmarks, even
exceeding GPT4 in the AlpacaEval benchmark, which evaluated models for
instruction following. We think that its ability to better emulate human
conversations will make it better at coming up with questions.

\textbf{Command-r:} Command-R is an LLM optimized for conversations with long
contexts. It was developed by Cohere Labs, to have grounded generation
capabilities. Essentially meaning that it is optimized for generating
responses based on supplied texts or documents. For this model we believe
that it might be a good guesser as it does well for long contexts and by
extension could perform well, even with several small questions.

\textbf{Cipher:} Cypher is a model that is optimized for human-like
conversations where the context/topic remains the same. It is good at
giving direct/precise answers and uses low memory. Its outputs are well
structured.

In 20 Questions, since the topic remains the same, Cypher might perform well as
it would be less likely to stray away from the topic when asking questions to
narrow the game down.

\textbf{Spark-X2.5-4B:} General purpose language model specialized for writing,
reasoning and conversations. A particular advantage of this model is that it
has a lower overhead compared to models which have a larger context. We think
that its specialization in reasoning will help it perform well in a game of 20
Questions.

\textbf{Thoxhuman-chat-8b:} Thox is a digital human conversation model that
can adapt its persona based on what is specified. It is good at emulating
natural dialogue and does a decent job when it comes to answering with
verified facts. It is based on the Quen3 model. Since Thox is good with facts
it might be able to guess quicker, before reaching question 20. Its natural
dialogue emulation may also help it with posing questions when it plays as a
questioner/player.

\textbf{apollo-astralis-4b:} Apollo-Astralis is built for deep logical
thinking and is customized for Human-AI collaboration. It has strong
conversational reasoning and communicates in a way that emulates human
empathy. It can work with RAM as low as 4 GB and occupy just about 2.5GB of
disk space. We think its specialization in logical thinking might make it good
for 20 Questions. Also it would be interesting to see if lightweight models can
achieve this.

\textbf{billy:} Billy stands for ``Better Interaction Language Learning Yield''
and it is customized for conversation tasks and sounding less bot like, and
more like a real human in a casual conversation. It only occupies 4.1 GB of
space and is relatively lightweight. We think that its customization for
conversational tasks will make it better at posing questions more naturally.

We may also use historical question sets available for 20 Questions games and
see how well each model performs with those legacy questions. This would give
us a good benchmark to compare each model.

There was a \textit{Twenty Questions} show that aired in 1949 on WOR-TV and
NBC. Although not all episodes are available (some records were lost) we think
it would be possible to use questions and topics from this show to create our
own benchmark to test models against.

\url{https://www.imdb.com/title/tt0320997/}

We also found the Clips 20Q Dataset from Hugging Face, which we could use to
make a similar benchmark.

\url{https://huggingface.co/datasets/clips/20Q/viewer/default/train}

\section{Evaluation Plan}

Our evaluation will focus on four dimensions: overall game performance, word
selection, question quality, and answer accuracy. Game performance will be
measured by win rate and the average number of questions needed to identify
the hidden word. We will test all pairings of at least four LLMs, with each
model serving as both questioner and answerer, to identify differences between
roles and self-play. Word selection will be evaluated by comparing the
distribution and diversity of model-selected words across categories (nouns,
places, people, and general vocabulary) against a uniform random baseline.
Question quality will be measured by how effectively questions reduce the
remaining candidate space, using information gain to assess how closely
questions divide possibilities in half. Answer accuracy will be evaluated by
comparing responses against known facts about the hidden word, measuring
truthfulness and consistency while identifying misleading or contradictory
answers.

We also plan to identify what characteristics make sentences/questions sound
more human. We will use this to come up with a metric that would help us
evaluate the humanness of each LLM.

\subsection*{Scoring Functions}

Generally we will have to develop the following scoring functions:

\begin{itemize}[leftmargin=*]
    \item Scoring functions for guess performance.
    \item Scoring function for how well the questions it asked were.
    \item Scoring on how varied the answers are.
    \item Scoring on how truthful the answerer is.
    \item Scoring on ``humanness'' of questions asked.
    \item Scoring on computational efficiency.
\end{itemize}

If time and resources permit, we will also evaluate computational efficiency
using metrics such as peak RAM and VRAM usage, response time, and token
consumption. This will allow us to examine the tradeoff between game
performance and resource requirements, identifying models that achieve
comparable results with fewer computational resources.

\section{Next Steps}

Currently, we have a rudimentary 20 questions game set up with 1 model
playing both the guesser and the answerer. As our next steps, we should allow
for multiple models to be loaded, and for multiple games to be run in sequence
to allow for better data collection. After collecting the data for multiple
models, and deciding on at least 4 to move on, we need to collect data for
each possible combination, and score each model based on the statistics we
chose earlier. If we have time later, we can try to run a larger model and see
how it will compare to the regular sized models we can run on our own hardware
by using Colab, or another available server. Additionally we could try to use
an API model to compare how the open-source vs proprietary models can tackle
this game.

Our code is located at
\url{https://github.com/n-man13/CSE519-Final-Project}.

\section{References}

\begin{enumerate}[leftmargin=*]
    \item KUMAR, CRS. ``Large Language Models (LLM) for Automating 20
    Questions Game.''
    \item Zoe Mongan, Luke Sernau, Will Lifferth, Bovard Doerschuk-Tiberi, Ryan
    Holbrook, Will Cukierski, Addison Howard. (2024). \textit{LLM 20
    Questions}. Kaggle.
    \url{https://kaggle.com/competitions/llm-20-questions}
    (Last accessed on 11/8/2026).
    \item Zhang, Yizhe, Jiarui Lu, and Navdeep Jaitly. ``Probing the
    multi-turn planning capabilities of LLMs via 20 question games.''
    \textit{Proceedings of the 62nd Annual Meeting of the Association for
    Computational Linguistics (Volume 1: Long Papers)}. 2024.
    \item \textit{20 questions (TV series 1949-- )} -- IMDb. (n.d.).
    \url{https://www.imdb.com/title/tt0320997/}
    \item CLiPS. (n.d.). \textit{20Q}. Hugging Face. Retrieved November 12,
    2026, from \url{https://huggingface.co/datasets/clips/20Q}
    \item Deepseek-ai/DeepSeek-R1-0528 --- Hugging Face. (2025, May 29).
    Hugging Face.
    \url{https://huggingface.co/deepseek-ai/DeepSeek-R1-0528}
    \item Moonshotai/Kimi-K2.6 --- Hugging Face. (2026, July 27). Hugging Face.
    \url{https://huggingface.co/moonshotai/Kimi-K2.6}
\end{enumerate}

\end{document}
